\chapter{Introducción}

La locomoción de robots cuadrúpedos ha sido ampliamente estudiada debido a su capacidad para desplazarse sobre terrenos irregulares y mantener la estabilidad ante perturbaciones. Plataformas como MIT Cheetah~3 y ANYmal muestran que una arquitectura mecánica articulada, combinada con estimación de estado y control de cuerpo completo, permite ejecutar desplazamientos robustos en condiciones experimentales exigentes~\cite{bledt2018cheetah3,hutter2016anymal}.

El desempeño locomotor depende de la configuración estructural, la posición del cuerpo, las trayectorias de los pies y las fuerzas de contacto. Por esta razón, el modelado cinemático y dinámico, el cálculo de las referencias articulares y la coordinación de las fases de apoyo y oscilación son elementos centrales del diseño de la marcha~\cite{zamani2011dynamics,bellicoso2018dynamic}.

Los enfoques modernos incorporan el estado del robot y la interacción pie--suelo para ajustar el movimiento. El control predictivo basado en modelo y la optimización de cuerpo completo permiten imponer restricciones de contacto y de los actuadores, además de responder a perturbaciones durante la locomoción~\cite{dicarlo2018mpc,neunert2018wholebody}. Estos resultados evidencian que el desarrollo de un sistema cuadrúpedo requiere integrar de manera coherente el diseño mecánico, el modelo matemático, la simulación y la arquitectura electrónica de control.

El presente proyecto desarrolló progresivamente el sistema de locomoción de la
plataforma Nova Spot Micro. Se construyeron modelos matemáticos y descripciones
para dos simuladores, se implementó una arquitectura de control en ROS~2 y se
obtuvieron líneas base reproducibles para postura, marcha paso y gateo. Además,
se incorporaron métricas de pose y articulaciones, contacto de los cuatro pies,
una IMU simulada, margen de estabilidad nominal y un supervisor de seguridad.

El desarrollo se organizó como una secuencia verificable: cada modificación de
la marcha se vinculó con su configuración, pruebas automatizadas, bolsa de datos
y análisis. Este criterio permitió identificar que el avance visual del robot
no garantizaba el cumplimiento del patrón de contactos previsto. Los resultados
presentados distinguen los componentes comprobados en simulación de las etapas
físicas y de aprendizaje por refuerzo que aún no han producido evidencia final.
